\chapter{PHÂN TÍCH THIẾT KẾ, TRIỂN KHAI VÀ ĐÁNH GIÁ HỆ THỐNG}

Agent:

Chương 3 đã cung cấp nền tảng lý thuyết và luận giải lựa chọn công nghệ cho toàn bộ hệ thống. Dựa trên nền tảng đó, chương 4 trình bày quá trình phân tích thiết kế, xây dựng và đánh giá hệ thống theo quy trình kỹ thuật phần mềm hoàn chỉnh. Năm nội dung chính được trình bày bao gồm: thiết kế kiến trúc ba tầng và tổng quan thành phần hệ thống trong mục 4.1, thiết kế chi tiết giao diện, lớp phần mềm và cơ sở dữ liệu trong mục 4.2, thống kê kết quả xây dựng và minh họa các chức năng cốt lõi trong mục 4.3, kế hoạch và kết quả kiểm thử hộp đen trong mục 4.4, và cấu hình triển khai trên môi trường phát triển lẫn production AWS trong mục 4.5.

\section{Thiết kế kiến trúc}

\subsection{Lựa chọn kiến trúc phần mềm}

Hệ thống áp dụng kiến trúc \textbf{ba tầng} (3-Tier Architecture) kết hợp với mô hình
\textbf{REST API / SPA}:

\begin{itemize}
  \item \textbf{Tầng trình bày (Presentation Tier):} Ứng dụng web đơn trang (SPA) React, chạy
    hoàn toàn trên trình duyệt, giao tiếp với backend qua HTTP REST~API.
  \item \textbf{Tầng logic nghiệp vụ (Business Logic Tier):} Django REST Framework xử lý toàn
    bộ logic dịch thuật, quản lý người dùng, glossary và lưu trữ.
  \item \textbf{Tầng dữ liệu (Data Tier):} PostgreSQL lưu dữ liệu có cấu trúc; AWS~S3 lưu
    tệp nhị phân; Google Cloud Storage lưu tệp CSV glossary.
\end{itemize}

Kiến trúc này được lựa chọn vì: (1)~tách biệt rõ ràng frontend và backend, dễ scale độc lập;
(2)~frontend SPA không cần server-side rendering nên phù hợp với giao diện nặng tương tác;
(3)~backend stateless (trừ DB) phù hợp triển khai sau Load Balancer.

\subsection{Thiết kế tổng quan}

Agent:

\begin{figure}[H]
\centering
\fbox{\parbox{0.9\textwidth}{\centering\vspace{1cm}
\textit{[Hình: Biểu đồ kiến trúc tổng quan]}\\[0.5cm]
% Mermaid source:
% flowchart TB
%     Browser["Trình duyệt\n(React SPA)"]
%     ALB["AWS ALB\n+ Cognito OIDC"]
%     BE["Django REST Framework\n(Backend API)"]
%     PG["PostgreSQL\n(Database)"]
%     S3["AWS S3\n(File Storage)"]
%     GCS["Google Cloud Storage\n(Glossary CSV)"]
%     GCTA["Google Cloud\nTranslation API v3"]
%     AZR["Azure Document\nIntelligence (OCR)"]
%     CVAPI["ConvertAPI\n(Format Conversion)"]
%
%     Browser -->|"HTTPS REST"| ALB
%     ALB -->|"x-amzn-oidc-data"| BE
%     BE -->|"ORM"| PG
%     BE -->|"boto3"| S3
%     BE -->|"translate_v3"| GCTA
%     GCTA -->|"glossary CSV"| GCS
%     BE -->|"OCR request"| AZR
%     BE -->|"convert request"| CVAPI
\textbf{Image Description:} Kiến trúc hệ thống gồm: Trình duyệt (React SPA) → AWS ALB (cân
bằng tải + Cognito OIDC) → Django REST Framework Backend. Backend kết nối PostgreSQL (ORM),
AWS S3 (boto3), Google Cloud Translation API v3, Google Cloud Storage (glossary CSV), Azure
Document Intelligence (OCR), ConvertAPI (chuyển đổi định dạng).
\vspace{0.5cm}}}
\caption{Biểu đồ kiến trúc tổng quan hệ thống}
\label{fig:architecture}
\end{figure}

\subsection{Thiết kế chi tiết các gói}

Agent:

\begin{figure}[H]
\centering
\fbox{\parbox{0.9\textwidth}{\centering\vspace{1cm}
\textit{[Hình: Biểu đồ UML Package]}\\[0.5cm]
\textbf{Image Description:} Biểu đồ package UML gồm 3 package chính:
\newline
(1) \textbf{Frontend} (React SPA): pages (HomePage, TextTranslation, TranslationResults,
FileHistory, CommonLibraryManagement, PrivateLibrary, MyProfile, AccountManagement,
KeywordStatsAdmin, LanguageManagement, FormatConversionPage, ConversionResults),
components (UploadFileNew, LanguageSelector, KeywordDetailModal, KeywordFormModal,
Notifications, LibraryActionButtons), contexts (AuthContext),
hooks (useCompanyLanguages, useLibraryLanguages),
services (translationService, keywordService, languageService, userService, notificationService).
\newline
(2) \textbf{Backend} (Django REST): views (TranslateFileView, TranslateTextView, alb\_auth,
keyword, translated\_file, notification, keyword\_stats, convert, user, language),
models (CustomUser, PasswordResetToken, Language, Company, CompanyLanguage,
TranslatedFile, PrivateKeyword, KeywordSuggestion, KeywordQueue,
LibraryQueueSettings, Notification, RefreshToken, LanguageAuditLog),
services (glossary\_service, upload\_to\_s3, convert\_api),
authentication (alb\_authentication),
Translate\_v2 (translate\_docx, translate\_xlsx, translate\_pptx, translate\_pdf\_ocr,
translate\_text, detect\_lang, xml\_process).
\newline
(3) \textbf{Data Layer}: PostgreSQL, AWS S3, Google Cloud Storage (glossary CSV).
\vspace{0.5cm}}}
\caption{Biểu đồ UML Package}
\label{fig:package}
\end{figure}

Agent:
\begin{lstlisting}[language=bash, caption={PlantUML source -- Biểu đồ Package (dán vào https://www.plantuml.com/plantuml)}]
@startuml
skinparam packageStyle rectangle

package "Frontend (React SPA)" {
  package "pages" {
    [HomePage] [TextTranslation] [TranslationResults]
    [FileHistory] [CommonLibraryManagement] [PrivateLibrary]
    [MyProfile] [AccountManagement] [KeywordStatsAdmin]
    [LanguageManagement] [FormatConversionPage]
  }
  package "components" {
    [UploadFileNew] [LanguageSelector] [KeywordDetailModal]
    [LibraryActionButtons] [Notifications]
  }
  package "context & hooks" {
    [AuthContext] [useCompanyLanguages] [useLibraryLanguages]
  }
  package "services-fe" {
    [translationService] [keywordService]
    [languageService] [userService] [notificationService]
  }
}

package "Backend (Django REST)" {
  package "views" {
    [TranslateFileView] [TranslateTextView]
    [keyword_views] [language_views] [user_views]
    [notification_views] [keyword_stats_views]
  }
  package "models" {
    [CustomUser] [Language] [Company] [CompanyLanguage]
    [TranslatedFile] [PrivateKeyword] [KeywordSuggestion]
    [KeywordQueue] [Notification] [RefreshToken] [LanguageAuditLog]
  }
  package "services-be" {
    [glossary_service] [upload_to_s3] [convert_api]
  }
  package "Translate_v2" {
    [translate_docx] [translate_xlsx] [translate_pptx]
    [translate_pdf_ocr] [translate_text] [detect_lang]
  }
}

package "Data Layer" {
  database "PostgreSQL"
  storage "AWS S3"
  storage "Google Cloud Storage"
}

[translationService] --> [TranslateFileView]
[keywordService] --> [keyword_views]
[languageService] --> [language_views]
[TranslateFileView] --> [translate_docx]
[TranslateFileView] --> [translate_xlsx]
[TranslateFileView] --> [translate_pdf_ocr]
[views] --> "PostgreSQL"
[upload_to_s3] --> "AWS S3"
[glossary_service] --> "Google Cloud Storage"
@enduml
\end{lstlisting}

\section{Thiết kế chi tiết}

\subsection{Thiết kế giao diện}

Hệ thống giao diện web bao gồm các trang chính sau:

\begin{enumerate}
  \item \textbf{Trang chủ / Upload tài liệu:} Widget tải tệp lên S3, chọn ngôn ngữ nguồn/đích
    (dropdown danh sách ngôn ngữ đã kích hoạt theo công ty, lấy động qua
    \texttt{useCompanyLanguages}), chọn chế độ thư viện, nút dịch và danh sách kết quả.
  \item \textbf{Trang dịch văn bản:} Textarea nhập văn bản, chọn ngôn ngữ, hiển thị kết quả
    dịch.
  \item \textbf{Trang lịch sử dịch:} Bảng danh sách TranslatedFile với phân trang, tìm kiếm,
    lọc theo ngôn ngữ.
  \item \textbf{Trang thư viện cá nhân:} Bảng PrivateKeyword với cột theo từng ngôn ngữ đang
    kích hoạt của công ty (dùng \texttt{translations}~JSONField), CRUD modal.
  \item \textbf{Trang thư viện chung (Library Keeper/Admin):} Quản lý KeywordSuggestion, duyệt
    đề xuất, tab phân loại (đang chờ / đã duyệt / đã từ chối).
  \item \textbf{Trang quản lý người dùng (Admin):} Bảng danh sách CustomUser, thay đổi vai trò.
  \item \textbf{Trang quản lý ngôn ngữ (Admin):} Bảng Language với tùy chọn kích hoạt/vô hiệu
    hóa ngôn ngữ theo công ty (\texttt{CompanyLanguage}), đặt ngôn ngữ mặc định, và xem nhật ký
    thay đổi (\texttt{LanguageAuditLog}).
  \item \textbf{Trang chuyển đổi định dạng:} Upload tệp PDF, chọn định dạng đích (DOCX/XLSX/PPTX),
    tải kết quả chuyển đổi về.
  \item \textbf{Trang thống kê từ khóa (Admin):} Biểu đồ recharts thống kê tần suất sử dụng
    từ khóa.
\end{enumerate}

Agent:

\begin{figure}[H]
\centering
\fbox{\parbox{0.9\textwidth}{\centering\vspace{2cm}
\textit{[Hình: Wireframe trang chủ Upload]}\\[0.5cm]
\textbf{Image Description:} Màn hình trang chủ gồm: thanh điều hướng trên cùng (logo, menu),
khu vực chính gồm widget kéo-thả tệp, dropdown chọn ngôn ngữ nguồn (auto-detect hoặc chọn
từ danh sách ngôn ngữ đã kích hoạt của công ty), danh sách checkbox chọn nhiều ngôn ngữ đích, radio chọn chế độ thư viện
(Không dùng / Thư viện chung / Thư viện cá nhân), nút ``Dịch tài liệu''. Phía dưới là khu vực
kết quả hiển thị danh sách các tệp đã dịch với nút tải về.
\vspace{1cm}}}
\caption{Wireframe trang Upload và dịch tài liệu}
\label{fig:ui_upload}
\end{figure}

\subsection{Thiết kế lớp}

\subsubsection{Lớp TranslateFileView}

\begin{figure}[H]
\centering
\fbox{\parbox{0.9\textwidth}{\centering\vspace{0.5cm}
\textit{[Hình: Class Diagram TranslateFileView]}\\[0.3cm]
\textbf{Image Description:} Class diagram PlantUML:
\newline
\texttt{TranslateFileView} kế thừa \texttt{APIView} (DRF), có thuộc tính
\texttt{permission\_classes = [IsAuthenticated]}. Phương thức: \texttt{post(request)} →
gọi \texttt{set\_translation\_context()}, vòng lặp qua \texttt{target\_languages}, gọi
\texttt{translate\_document()}, tạo \texttt{TranslatedFile}, gọi
\texttt{clear\_translation\_context()}.
\newline
\texttt{translate\_document(file\_path, target\_language, ...)} → \texttt{is\_pdf\_truly\_editable()}
→ phân nhánh: editable → \texttt{convertapi\_pdf\_to\_docx()} + \texttt{translate\_docx()};
scan → \texttt{translate\_pdf\_ocr()}; docx → \texttt{translate\_docx()}; xlsx →
\texttt{translate\_xlsx()}; pptx → \texttt{translate\_pptx()}.
\vspace{0.5cm}}}
\caption{Class diagram -- TranslateFileView và translate\_document}
\label{fig:class_translate}
\end{figure}

Agent:
\begin{lstlisting}[language=bash, caption={PlantUML source -- Class Diagram TranslateFileView}]
@startuml
class TranslateFileView {
  +permission_classes: [IsAuthenticated]
  +post(request: Request): Response
  -_download_from_s3(url): str
  -_upload_to_s3(path, lang): str
}

class translate_document {
  +file_path: str
  +target_lang: str
  +source_lang: str
  +library_mode: str
  +is_pdf_truly_editable(pdf_path: str): bool
}

class set_translation_context {
  +library_mode: str
  +user_id: int
}

class clear_translation_context

class translate_docx << (F,lightblue) function >>
class translate_xlsx << (F,lightblue) function >>
class translate_pptx << (F,lightblue) function >>
class translate_pdf_ocr << (F,lightblue) function >>
class convertapi_pdf_to_docx << (F,lightblue) function >>

TranslateFileView --> set_translation_context : "1. set context"
TranslateFileView --> translate_document : "2. per target lang"
TranslateFileView --> clear_translation_context : "3. finally"
translate_document --> translate_docx : "docx / pdf-editable"
translate_document --> translate_xlsx : "xlsx / xls / csv"
translate_document --> translate_pptx : "pptx"
translate_document --> translate_pdf_ocr : "pdf scan"
translate_document --> convertapi_pdf_to_docx : "pdf-editable pre-step"
@enduml
\end{lstlisting}

\subsubsection{Lớp dữ liệu Model}

\begin{figure}[H]
\centering
\fbox{\parbox{0.9\textwidth}{\centering\vspace{0.5cm}
\textit{[Hình: Class Diagram các Model]}\\[0.3cm]
\textbf{Image Description:} Class diagram 13 Django Model:
\newline
\textbf{CustomUser}: email (unique), first\_name, last\_name, department,
role (User/Admin/Library Keeper), is\_active, FK(Company null), FK(Language preferred null), date\_joined.
\newline
\textbf{Language}: code (unique, e.g.\ ``vi'', ``ja'', ``zh-CN''), name, native\_name, flag\_emoji, is\_active, sort\_order.
\newline
\textbf{Company}: name (unique), slug (unique), is\_active, created\_at.
\newline
\textbf{CompanyLanguage}: FK(Company), FK(Language), is\_enabled, is\_default, enabled\_at.
UNIQUE(company, language). Partial unique index on is\_default (migration 0032).
\newline
\textbf{TranslatedFile}: FK(CustomUser), original\_file\_url, original\_file\_name,
translated\_file\_url, original\_language, target\_language, file\_type, created\_at.
\newline
\textbf{PrivateKeyword}: FK(CustomUser), translations (JSONField, default=dict),
note, FK(KeywordSuggestion null), created\_at, updated\_at.
\newline
\textbf{KeywordSuggestion}: FK(CustomUser), translations (JSONField), suggestion\_count,
frequency\_percentage, status (pending/approved/rejected), FK(approved\_by null), created\_at, updated\_at.
\newline
\textbf{KeywordQueue}: FK(CustomUser), translations (JSONField), is\_processed (db\_index), processed\_at, created\_at.
\newline
\textbf{Notification}: FK(CustomUser), title, message, read, details, keyword\_details (JSONField), created\_at.
\newline
\textbf{RefreshToken}: FK(CustomUser), token\_hash (SHA-256, unique, db\_index), expires\_at, created\_at.
\newline
\textbf{LanguageAuditLog}: FK(actor=CustomUser), FK(Company), FK(Language),
action (enable/disable/set\_default/add\_language), note, created\_at.
\newline
\textbf{LibraryQueueSettings}: pk=1 (singleton), min\_suggesters\_for\_queue, updated\_at.
\newline
\textbf{PasswordResetToken}: email, token (OTP 6 chữ số), created\_at, expires\_at (TTL 10 phút).
\vspace{0.5cm}}}
\caption{Class diagram -- Django Models}
\label{fig:class_models}
\end{figure}

Agent:
\begin{lstlisting}[language=bash, caption={PlantUML source -- Class Diagram Django Models}]
@startuml
class CustomUser {
  +email: EmailField (unique)
  +first_name: CharField
  +last_name: CharField
  +department: CharField
  +role: CharField [User|Admin|Library Keeper]
  +is_active: BooleanField
  +company: FK(Company, null)
  +preferred_language: FK(Language, null)
  +date_joined: DateTimeField
}
class Language {
  +code: CharField (unique)  -- vi, ja, zh-CN ...
  +name: CharField
  +native_name: CharField
  +flag_emoji: CharField
  +is_active: BooleanField
  +sort_order: PositiveSmallIntegerField
}
class Company {
  +name: CharField (unique)
  +slug: SlugField (unique)
  +is_active: BooleanField
}
class CompanyLanguage {
  +company: FK(Company)
  +language: FK(Language)
  +is_enabled: BooleanField
  +is_default: BooleanField
  -- UNIQUE(company, language)
}
class TranslatedFile {
  +user: FK(CustomUser)
  +original_file_url: URLField
  +original_file_name: CharField
  +translated_file_url: URLField
  +original_language: CharField
  +target_language: CharField
  +file_type: CharField
}
class PrivateKeyword {
  +user: FK(CustomUser)
  +translations: JSONField  -- {en:..,ja:..,vi:..}
  +note: TextField
  +suggestion: FK(KeywordSuggestion, null)
}
class KeywordSuggestion {
  +user: FK(CustomUser)
  +translations: JSONField
  +suggestion_count: IntegerField
  +status: CharField [pending|approved|rejected]
  +approved_by: FK(CustomUser, null)
}
class KeywordQueue {
  +user: FK(CustomUser)
  +translations: JSONField
  +is_processed: BooleanField (db_index)
  +processed_at: DateTimeField
}
class Notification {
  +user: FK(CustomUser)
  +title: CharField
  +message: TextField
  +read: BooleanField
  +keyword_details: JSONField
}
class RefreshToken {
  +user: FK(CustomUser)
  +token_hash: CharField (SHA-256, unique)
  +expires_at: DateTimeField
}
class LanguageAuditLog {
  +actor: FK(CustomUser)
  +company: FK(Company)
  +language: FK(Language)
  +action: CharField [enable|disable|set_default|add_language]
  +note: TextField
}
class LibraryQueueSettings {
  +pk: 1 (singleton)
  +min_suggesters_for_queue: PositiveIntegerField
}

CustomUser "N" --> "1" Company
CustomUser "N" --> "0..1" Language : preferred
Company "1" --> "N" CompanyLanguage
Language "1" --> "N" CompanyLanguage
CustomUser "1" --> "N" TranslatedFile
CustomUser "1" --> "N" PrivateKeyword
CustomUser "1" --> "N" KeywordQueue
CustomUser "1" --> "N" KeywordSuggestion
KeywordSuggestion "1" --> "N" PrivateKeyword : linked
CustomUser "1" --> "N" RefreshToken
CustomUser "1" --> "N" Notification
@enduml
\end{lstlisting}

\subsubsection{Lớp GlossaryService}

\begin{figure}[H]
\centering
\fbox{\parbox{0.9\textwidth}{\centering\vspace{0.5cm}
\textit{[Hình: Class Diagram GlossaryService]}\\[0.3cm]
\textbf{Image Description:} Class diagram:
\newline
\texttt{GlossaryService (module)} gồm các hàm:
\texttt{create\_pair\_csv\_file(user, source\_lang, target\_lang)} → truy vấn
\texttt{KeywordSuggestion(approved)} + \texttt{PrivateKeyword(user)} → build pair rows →
ghi CSV tạm.
\texttt{upload\_csv\_to\_gcs(csv\_path)} → \texttt{storage.Client().bucket().blob().upload()}.
\texttt{manage\_glossary(glossary\_id, src, tgt, mode)} → gọi
\texttt{translate.TranslationServiceClient().create\_glossary()} hoặc
\texttt{update\_glossary()}.
\texttt{async\_manage\_common\_glossaries()} → background thread.
\texttt{async\_manage\_user\_glossaries(user)} → background thread cho private glossary.
\vspace{0.5cm}}}
\caption{Class diagram -- GlossaryService}
\label{fig:class_glossary}
\end{figure}

Agent:
\begin{lstlisting}[language=bash, caption={PlantUML source -- Class Diagram GlossaryService}]
@startuml
package "glossary_service (module)" {
  class create_pair_csv_file << (F,lightblue) >> {
    +user: CustomUser
    +source_lang: str
    +target_lang: str
    -- queries KeywordSuggestion(approved) --
    -- builds pair rows -> writes CSV temp --
    +return: csv_path: str
  }
  class upload_csv_to_gcs << (F,lightblue) >> {
    +csv_path: str
    -- storage.Client().bucket().blob().upload() --
    +return: gcs_uri: str
  }
  class manage_glossary << (F,lightblue) >> {
    +glossary_id: str
    +src: str
    +tgt: str
    +mode: int  -- 0=create, 1=update
    -- TranslationServiceClient().create_glossary() --
    -- or update_glossary() --
  }
  class async_manage_common_glossaries << (F,lightblue) >> {
    -- background Thread --
    -- iterates 45 language pairs (C(10,2) via generate_pairs_from_db()) --
    -- calls create_pair_csv_file + upload + manage_glossary --
  }
  class async_manage_user_glossaries << (F,lightblue) >> {
    +user: CustomUser
    -- background Thread --
    -- private glossary: toray_glossary_A_B_user_ID --
  }
}

create_pair_csv_file --> upload_csv_to_gcs
upload_csv_to_gcs --> manage_glossary
async_manage_common_glossaries --> create_pair_csv_file
async_manage_user_glossaries --> create_pair_csv_file
@enduml
\end{lstlisting}

\subsection{Thiết kế cơ sở dữ liệu}

Hệ thống sử dụng 12 bảng chính trong PostgreSQL:

Agent:

\begin{figure}[H]
\centering
\fbox{\parbox{0.9\textwidth}{\centering\vspace{0.5cm}
\textit{[Hình: Biểu đồ ERD]}\\[0.3cm]
% Mermaid ERD source:
% erDiagram
%     CUSTOM_USER {
%         int id PK
%         varchar email UK
%         varchar first_name
%         varchar last_name
%         varchar department
%         varchar role
%         bool is_active
%         varchar cognito_sub UK
%         varchar auth_provider
%         datetime date_joined
%     }
%     TRANSLATED_FILE {
%         int id PK
%         int user_id FK
%         varchar original_file_url
%         varchar original_file_name
%         varchar translated_file_url
%         varchar original_language
%         varchar target_language
%         varchar file_type
%         datetime created_at
%     }
%     PRIVATE_KEYWORD {
%         int id PK
%         int user_id FK
%         int suggestion_id FK
%         text japanese
%         text english
%         text vietnamese
%         text chinese_traditional
%         text chinese_simplified
%         text thai
%         text bengali
%         text hindi
%         text indonesian
%         text oriya
%         text note
%     }
%     KEYWORD_SUGGESTION {
%         int id PK
%         int user_id FK
%         int approved_by_id FK
%         text japanese
%         text english
%         text vietnamese
%         text chinese_traditional
%         text chinese_simplified
%         text thai
%         text bengali
%         text hindi
%         text indonesian
%         text oriya
%         int suggestion_count
%         float frequency_percentage
%         varchar status
%     }
%     KEYWORD_QUEUE {
%         int id PK
%         int user_id FK
%         [10 language fields]
%         bool is_processed
%         datetime processed_at
%         datetime created_at
%     }
%     NOTIFICATION {
%         int id PK
%         int user_id FK
%         varchar title
%         text message
%         bool read
%         text details
%         json keyword_details
%         datetime created_at
%     }
%     LIBRARY_QUEUE_SETTINGS {
%         int id PK
%         int min_suggesters_for_queue
%         datetime updated_at
%     }
%
%     CUSTOM_USER ||--o{ TRANSLATED_FILE : "dịch"
%     CUSTOM_USER ||--o{ PRIVATE_KEYWORD : "sở hữu"
%     CUSTOM_USER ||--o{ KEYWORD_QUEUE : "đề xuất"
%     CUSTOM_USER ||--o{ KEYWORD_SUGGESTION : "tạo"
%     CUSTOM_USER ||--o{ NOTIFICATION : "nhận"
%     KEYWORD_SUGGESTION ||--o{ PRIVATE_KEYWORD : "liên kết"
%     CUSTOM_USER ||--o{ KEYWORD_SUGGESTION : "duyệt (approved_by)"
\textbf{Image Description:} Biểu đồ ERD với 12 thực thể:
CUSTOM\_USER (id, email unique, first\_name, last\_name, department, role, FK company null, FK preferred\_language null, is\_active, date\_joined);
LANGUAGE (id, code unique, name, native\_name, flag\_emoji, is\_active, sort\_order);
COMPANY (id, name unique, slug unique, is\_active, created\_at);
COMPANY\_LANGUAGE (id, FK company, FK language, is\_enabled, is\_default, enabled\_at --- UNIQUE company+language);
TRANSLATED\_FILE (id, FK user, original\_file\_url, original\_file\_name, translated\_file\_url, original\_language, target\_language, file\_type, created\_at);
PRIVATE\_KEYWORD (id, FK user, FK suggestion null, translations JSONField, note, created\_at, updated\_at);
KEYWORD\_SUGGESTION (id, FK user, FK approved\_by null, translations JSONField, suggestion\_count, frequency\_percentage, status, created\_at, updated\_at);
KEYWORD\_QUEUE (id, FK user, translations JSONField, is\_processed, processed\_at, created\_at);
NOTIFICATION (id, FK user, title, message, read, details, keyword\_details JSONField, created\_at);
REFRESH\_TOKEN (id, FK user, token\_hash SHA-256 unique, expires\_at, created\_at);
LANGUAGE\_AUDIT\_LOG (id, FK actor, FK company, FK language, action, note, created\_at);
LIBRARY\_QUEUE\_SETTINGS (pk=1 singleton, min\_suggesters\_for\_queue, updated\_at).
Quan hệ chính: CUSTOM\_USER N-1 COMPANY; CUSTOM\_USER N-1 LANGUAGE (preferred);
COMPANY N-N LANGUAGE qua COMPANY\_LANGUAGE; CUSTOM\_USER 1-N TRANSLATED\_FILE;
CUSTOM\_USER 1-N PRIVATE\_KEYWORD; CUSTOM\_USER 1-N KEYWORD\_QUEUE;
CUSTOM\_USER 1-N KEYWORD\_SUGGESTION; KEYWORD\_SUGGESTION 1-N PRIVATE\_KEYWORD;
CUSTOM\_USER 1-N REFRESH\_TOKEN; CUSTOM\_USER 1-N LANGUAGE\_AUDIT\_LOG.
\vspace{0.5cm}}}
\caption{Biểu đồ ERD cơ sở dữ liệu}
\label{fig:erd}
\end{figure}

Agent:
\begin{lstlisting}[language=bash, caption={Mermaid ERD source -- 12 entities (dán vào https://mermaid.live)}]
erDiagram
    CUSTOM_USER {
        int id PK
        varchar email UK
        varchar first_name
        varchar last_name
        varchar department
        varchar role
        bool is_active
        int company_id FK
        int preferred_language_id FK
        datetime date_joined
    }
    LANGUAGE {
        int id PK
        varchar code UK
        varchar name
        varchar native_name
        varchar flag_emoji
        bool is_active
        int sort_order
    }
    COMPANY {
        int id PK
        varchar name UK
        varchar slug UK
        bool is_active
    }
    COMPANY_LANGUAGE {
        int id PK
        int company_id FK
        int language_id FK
        bool is_enabled
        bool is_default
        datetime enabled_at
    }
    TRANSLATED_FILE {
        int id PK
        int user_id FK
        varchar original_file_url
        varchar original_file_name
        varchar translated_file_url
        varchar original_language
        varchar target_language
        varchar file_type
        datetime created_at
    }
    PRIVATE_KEYWORD {
        int id PK
        int user_id FK
        int suggestion_id FK
        json translations
        text note
        datetime created_at
        datetime updated_at
    }
    KEYWORD_SUGGESTION {
        int id PK
        int user_id FK
        int approved_by_id FK
        json translations
        int suggestion_count
        float frequency_percentage
        varchar status
        datetime created_at
    }
    KEYWORD_QUEUE {
        int id PK
        int user_id FK
        json translations
        bool is_processed
        datetime processed_at
        datetime created_at
    }
    NOTIFICATION {
        int id PK
        int user_id FK
        varchar title
        text message
        bool read
        json keyword_details
        datetime created_at
    }
    REFRESH_TOKEN {
        int id PK
        int user_id FK
        varchar token_hash UK
        datetime expires_at
        datetime created_at
    }
    LANGUAGE_AUDIT_LOG {
        int id PK
        int actor_id FK
        int company_id FK
        int language_id FK
        varchar action
        text note
        datetime created_at
    }
    LIBRARY_QUEUE_SETTINGS {
        int id PK
        int min_suggesters_for_queue
        datetime updated_at
    }

    CUSTOM_USER }o--|| COMPANY : "belongs to"
    CUSTOM_USER }o--o| LANGUAGE : "prefers"
    COMPANY ||--o{ COMPANY_LANGUAGE : "activates"
    LANGUAGE ||--o{ COMPANY_LANGUAGE : "activated in"
    CUSTOM_USER ||--o{ TRANSLATED_FILE : "translates"
    CUSTOM_USER ||--o{ PRIVATE_KEYWORD : "owns"
    CUSTOM_USER ||--o{ KEYWORD_QUEUE : "queues"
    CUSTOM_USER ||--o{ KEYWORD_SUGGESTION : "suggests"
    CUSTOM_USER ||--o{ NOTIFICATION : "receives"
    CUSTOM_USER ||--o{ REFRESH_TOKEN : "holds"
    KEYWORD_SUGGESTION ||--o{ PRIVATE_KEYWORD : "linked to"
\end{lstlisting}

\section{Xây dựng ứng dụng}

\subsection{Thư viện và công cụ sử dụng}

\begin{table}[H]
\centering
\caption{Thư viện và công cụ chính được sử dụng}
\label{tab:tools}
\renewcommand{\arraystretch}{1.3}
\begin{tabularx}{\textwidth}{|l|l|X|}
\hline
\textbf{Công cụ / Thư viện} & \textbf{Phiên bản} & \textbf{Mục đích} \\
\hline
Visual Studio Code     & --    & IDE phát triển \\
\hline
Git                    & --    & Quản lý phiên bản \\
\hline
Docker                 & --    & Container hóa ứng dụng \\
\hline
React + Vite           & 18/6  & Frontend SPA \\
\hline
Django + DRF           & 6/3.16& REST API backend \\
\hline
PostgreSQL             & 12+   & Cơ sở dữ liệu production \\
\hline
google-cloud-translate & v3    & Dịch thuật với glossary \\
\hline
azure-ai-documentintelligence & -- & OCR tài liệu scan \\
\hline
convertapi             & --    & Chuyển đổi PDF→DOCX, XLS→XLSX \\
\hline
python-docx / lxml     & --    & Xử lý DOCX trực tiếp XML \\
\hline
openpyxl               & --    & Xử lý XLSX cell-by-cell \\
\hline
python-pptx            & --    & Xử lý PPTX shapes \\
\hline
pypdf                  & --    & Phân tích cấu trúc PDF \\
\hline
boto3                  & --    & Upload/download S3 \\
\hline
\end{tabularx}
\end{table}

\subsection{Kết quả đạt được}

\begin{table}[H]
\centering
\caption{Thống kê mã nguồn}
\label{tab:codestats}
\renewcommand{\arraystretch}{1.3}
\begin{tabularx}{\textwidth}{|l|X|}
\hline
\textbf{Thành phần}     & \textbf{Số lượng} \\
\hline
Backend views           & 12 view files \\
\hline
Backend models          & 10 model files (13 Django model class) \\
\hline
Backend services        & 3 service modules (glossary, S3, ConvertAPI) \\
\hline
Backend migrations      & 35 migration files \\
\hline
Frontend pages (routes) & 15 routes \\
\hline
Frontend components     & 34 component files \\
\hline
Translation modules     & 6 format modules (docx, xlsx, pptx, pdf, pdf\_ocr, text) \\
\hline
Glossary pairs          & 45 cặp ngôn ngữ (C(10,2), \texttt{generate\_pairs\_from\_db()}) \\
\hline
Ngôn ngữ hỗ trợ         & 10 ngôn ngữ (seed ban đầu; cấu hình động qua LanguageManagement) \\
\hline
Định dạng tài liệu      & 6 (DOCX, XLSX, PPTX, PDF, XLS/XLSM/XLSB, CSV) \\
\hline
\end{tabularx}
\end{table}

\subsection{Minh họa các chức năng chính}

\subsubsection{Luồng dịch tài liệu DOCX}

Dưới đây là đoạn mã minh họa hàm \texttt{translate\_all\_xml\_in\_folder} trong module
\texttt{translate\_docx\_traditional.py} --- thực hiện dịch batch toàn bộ chuỗi văn bản trong
tệp DOCX bằng cách parse XML trực tiếp:

\begin{lstlisting}[language=Python, caption={Batch translation trong DOCX (translate\_docx\_traditional.py)}]
def translate_all_xml_in_folder(folder_path, glossary_id,
                                 source_lang, target_lang):
    all_xml_tasks = []
    texts_to_translate = set()

    for root_dir, dirs, files in os.walk(folder_path):
        for file in files:
            if not file.endswith(".xml"):
                continue
            xml_file_path = os.path.join(root_dir, file)
            if not _is_docx_content_xml(folder_path, xml_file_path):
                continue
            clean_and_merge_runs(xml_file_path)  # merge runs cung format
            # Thu thap text elements
            tree = ET.parse(xml_file_path, ET.XMLParser(...))
            for elem in tree.getroot().iter():
                if elem.tag.endswith("}t") and elem.text.strip():
                    texts_to_translate.add(elem.text.strip())

    # Dich BATCH toan bo chuoi duy nhat
    unique_texts = list(texts_to_translate)
    translated_list = translate_texts_batch(
        unique_texts, glossary_id, source_lang, target_lang)
    translation_map = dict(zip(unique_texts, translated_list))

    # Cap nhat va luu lai cac file XML
    for xml_file_path, tree, elements in all_xml_tasks:
        for elem in elements:
            if elem.text.strip() in translation_map:
                elem.text = html.unescape(
                    translation_map[elem.text.strip()])
        tree.write(xml_file_path, encoding="UTF-8")
\end{lstlisting}

\subsubsection{Phân loại PDF}

Hàm \texttt{is\_pdf\_truly\_editable()} phân tích từng trang PDF để phân loại:

\begin{lstlisting}[language=Python, caption={Phân loại PDF native-text vs scan (translate.py)}]
def is_pdf_truly_editable(pdf_path: str) -> bool:
    with open(pdf_path, "rb") as f:
        reader = PyPDF2.PdfReader(f)
        for i, page in enumerate(reader.pages, start=1):
            text = (page.extract_text() or "").strip()
            has_text = len(text) >= MIN_TEXT_CHARS       # 20 chars
            has_full_image = _page_full_page_image(page) # XObject ratio
            if has_text and not has_full_image:
                return True  # native text page found
    return False  # all pages are scan/image
\end{lstlisting}

\section{Kiểm thử}

\subsection{Kế hoạch kiểm thử}

Hệ thống áp dụng phương pháp kiểm thử hộp đen (Black-box testing), kiểm tra từng chức năng
dựa trên yêu cầu đặc tả. Bảng~\ref{tab:testcases} trình bày các ca kiểm thử chính.

\begin{table}[H]
\centering
\caption{Bảng ca kiểm thử}
\label{tab:testcases}
\renewcommand{\arraystretch}{1.3}
\begin{tabularx}{\textwidth}{|l|l|X|X|l|}
\hline
\textbf{ID} & \textbf{Nhóm} & \textbf{Mô tả} & \textbf{Kết quả mong đợi} & \textbf{KQ} \\
\hline
TC-01 & Xác thực & Đăng nhập email/mật khẩu hợp lệ & JWT trả về, chuyển trang & Đạt \\
\hline
TC-02 & Xác thực & Đăng nhập sai mật khẩu & Lỗi 401 & Đạt \\
\hline
TC-03 & Dịch tài liệu & Dịch DOCX sang tiếng Nhật & Tệp DOCX dịch trả về, định dạng giữ nguyên & Đạt \\
\hline
TC-04 & Dịch tài liệu & Dịch XLSX sang 3 ngôn ngữ cùng lúc & 3 tệp XLSX kết quả & Đạt \\
\hline
TC-05 & Dịch tài liệu & Dịch PDF native-text & Tệp DOCX dịch trả về & Đạt \\
\hline
TC-06 & Dịch tài liệu & Dịch PDF scan & Tệp PDF dịch trả về (qua OCR) & Đạt \\
\hline
TC-07 & Dịch tài liệu & Dịch PPTX sang tiếng Trung & Tệp PPTX dịch trả về & Đạt \\
\hline
TC-08 & Dịch văn bản & Dịch văn bản từ tiếng Nhật sang tiếng Việt & Văn bản dịch hiển thị & Đạt \\
\hline
TC-09 & Dịch văn bản & Dịch với thuật ngữ trong glossary & Thuật ngữ được dịch đúng theo glossary & Đạt \\
\hline
TC-10 & Thư viện & Thêm từ khóa vào Private Library & Bản ghi PrivateKeyword được tạo & Đạt \\
\hline
TC-11 & Thư viện & Đề xuất từ khóa lên Common Library & Bản ghi KeywordQueue được tạo & Đạt \\
\hline
TC-12 & Thư viện & Duyệt đề xuất (Library Keeper) & Status = approved; GCS cập nhật & Đạt \\
\hline
TC-13 & Thư viện & Từ chối đề xuất & Status = rejected; thông báo gửi user & Đạt \\
\hline
TC-14 & Quản trị & Admin thay đổi vai trò người dùng & Role được cập nhật trong DB & Đạt \\
\hline
TC-15 & Quản trị & Admin xem thống kê từ khóa & Biểu đồ recharts hiển thị đúng & Đạt \\
\hline
TC-16 & Lịch sử & Xem danh sách tệp đã dịch & Bảng TranslatedFile hiển thị đúng & Đạt \\
\hline
TC-17 & Phi chức năng & Dịch tệp DOCX 10 trang & Hoàn thành trong $\leq 30$s & Đạt \\
\hline
TC-18 & Phi chức năng & Truy cập endpoint không xác thực & Lỗi 401 & Đạt \\
\hline
\end{tabularx}
\end{table}

\section{Triển khai}

\subsection{Môi trường phát triển}

\begin{itemize}
  \item Backend: \texttt{python manage.py runserver} tại \texttt{http://127.0.0.1:8000}
  \item Frontend: \texttt{npm run dev} (Vite) tại \texttt{http://localhost:5173}
  \item Database: SQLite (fallback) hoặc PostgreSQL local
  \item Auth: JWT simplejwt (local login, không cần Cognito)
\end{itemize}

\subsection{Môi trường production}

\begin{figure}[H]
\centering
\fbox{\parbox{0.9\textwidth}{\centering\vspace{0.5cm}
\textit{[Hình: Sơ đồ triển khai production]}\\[0.3cm]
\textbf{Image Description:} Sơ đồ triển khai: Người dùng → HTTPS → AWS ALB (với Cognito OIDC
authentication action) → Target Group EC2/ECS chạy Docker containers:
container-frontend (Nginx phục vụ React build) và container-backend (Django + Gunicorn).
Backend kết nối tới: RDS PostgreSQL, AWS S3 (boto3), Google Cloud Translation API (HTTPS),
Azure Document Intelligence (HTTPS), ConvertAPI (HTTPS). Biến môi trường được inject qua
.env file hoặc AWS Secrets Manager.
\vspace{0.5cm}}}
\caption{Sơ đồ triển khai production trên AWS}
\label{fig:deployment}
\end{figure}

Agent:
\begin{lstlisting}[language=bash, caption={Mermaid source -- Sơ đồ triển khai production (dán vào https://mermaid.live)}]
graph TB
    Browser["Trình duyệt (HTTPS)"]
    Browser --> ALB["AWS Application\nLoad Balancer\n(Cognito OIDC auth action)"]

    ALB --> Nginx["Nginx Container\n(ports 80/443)"]

    subgraph Docker["Docker Compose — EC2/ECS"]
        Nginx --> FE["Frontend Container\n(Nginx + React build\nport 80 internal)"]
        Nginx --> BE["Backend Container\n(Django + Gunicorn\nport 8000 internal)"]
    end

    BE --> PG[("RDS PostgreSQL")]
    BE --> S3[("AWS S3\noriginal + translated files")]
    BE --> GCT["Google Cloud\nTranslation API v3\n(45 glossary pairs, dynamic)"]
    GCT --> GCS[("Google Cloud Storage\nCSV glossary files")]
    BE --> AZR["Azure Document\nIntelligence (OCR)"]
    BE --> CAPI["ConvertAPI\n(PDF→DOCX)"]
    BE --> SM["AWS Secrets Manager\n(env variables)"]
\end{lstlisting}

Hệ thống production được triển khai trên AWS với:
\begin{itemize}
  \item \textbf{Xác thực:} AWS ALB tích hợp Amazon Cognito OIDC, inject header
    \texttt{x-amzn-oidc-data} (JWT) vào mọi request. Backend middleware
    \texttt{alb\_authentication.py} giải mã header này, tra cứu hoặc tạo \texttt{CustomUser}
    dựa trên \texttt{cognito\_sub}.
  \item \textbf{Lưu trữ tệp:} AWS~S3 bucket lưu tệp gốc (\texttt{uploads/}) và tệp dịch
    (\texttt{translated/\{uuid\}.\{ext\}}) với tên UUID ngẫu nhiên.
  \item \textbf{Container:} Dockerfile riêng cho backend (Python + pip) và frontend
    (Node build + Nginx serve). Frontend Nginx cấu hình \texttt{try\_files} cho SPA routing.
\end{itemize}

Agent:

Kết chương 4. Chương 4 đã trình bày đầy đủ toàn bộ quá trình phân tích, thiết kế, xây dựng và đánh giá hệ thống. Kiến trúc ba tầng --- SPA React, Django REST Framework và lớp dữ liệu PostgreSQL/S3/GCS --- đảm bảo tính tách biệt rõ ràng và khả năng mở rộng độc lập của từng tầng. Cơ sở dữ liệu 12 bảng được thiết kế bao phủ toàn bộ nghiệp vụ: quản lý người dùng và công ty, cấu hình ngôn ngữ theo công ty (Language/Company/CompanyLanguage), tệp dịch, thuật ngữ ba cấp (PrivateKeyword/KeywordQueue/KeywordSuggestion), thông báo, token xác thực và nhật ký kiểm tra ngôn ngữ, trong khi lớp GlossaryService chịu trách nhiệm đồng bộ 45 cặp glossary lên Google Translation API tự động qua background thread (C(10,2)~=~45 với 10 ngôn ngữ mặc định, tự động mở rộng theo \texttt{generate\_pairs\_from\_db()}). Kết quả kiểm thử 18 ca kiểm thử hộp đen bao phủ toàn bộ nhóm chức năng --- xác thực, dịch tài liệu, dịch văn bản, quản lý thuật ngữ và quản trị hệ thống --- đều đạt yêu cầu đề ra. Hệ thống được triển khai production trên AWS với Cognito SSO, ALB và Docker container hóa, đáp ứng các yêu cầu phi chức năng về bảo mật và khả năng mở rộng. Chương tiếp theo sẽ tách riêng và phân tích chuyên sâu ba đóng góp kỹ thuật nổi bật nhất của đồ án trong quá trình xây dựng hệ thống.
